%! Function Model Selection and Construction — Unit 1, Lesson 1.7 — Lesson Plan
\documentclass[10pt]{article}
\usepackage{precalculus-article}
\usepackage{precalculus-boxes}
\usepackage{graphicx}
\graphicspath{{images/}}
\newcommand{\TallMath}[1]{$\displaystyle #1\rule[-1.4em]{0pt}{3.2em}$}
\newcommand{\CourseName}{Precalculus}
\newcommand{\MeetingLength}{60 minutes}
\newcommand{\UnitNumberName}{Unit 1: Functions and Change \quad}
\newcommand{\LessonNumberName}{Lesson 1.7: Function Model Selection and Construction}

\begin{document}
\begin{center}
    {\Large\bfseries \CourseName \\
    {\small\bfseries \UnitNumberName \LessonNumberName}}
\end{center}

\begin{tcolorbox}[colback=lilac, colframe=plum, boxrule=0.9pt, arc=2mm,
    left=2mm, right=2mm, top=1.5mm, bottom=1.5mm]
    \textbf{Primary Objective:} Given a situation rather than a function,
    students will \textbf{name which quantity is the question}, choose a model
    type for it --- from differences in a table or from shape clues in the
    context --- \textbf{build} the model by letting a constraint force the other
    quantities to answer to that question, and state the domain, the units, and
    the assumption the model rests on.
    \hfill\textit{\small (CED Topics 1.13 \& 1.14 --- Skills 2.A, 3.C, 1.C, 3.B)}
\end{tcolorbox}

\renewcommand{\arraystretch}{1.6}
\begin{skillbox}[Priority Ideas \& Skills]{goldbox}
\begin{tabularx}{\linewidth}{@{}X X@{}}
\textbf{AP Skills} &
\textbf{Key Understandings} \\
\midrule
\textbf{2.A} --- Identify information from graphical, numerical, analytical, and
verbal representations to answer a question or construct a model. \newline
\textbf{3.C} --- Support conclusions or choices with a logical rationale. &
\begin{itemize}[leftmargin=1.2em,itemsep=2pt,topsep=0pt]
  \item \textbf{The spine.} A situation has several quantities; a function
        accepts one question. Modeling begins by deciding which quantity is
        the \textbf{question}.
  \item A table has already decided --- the top row is the question. Over
        \textbf{equally spaced} inputs: constant first differences $\Rightarrow$
        linear; constant second $\Rightarrow$ quadratic; constant $n$th
        $\Rightarrow$ degree $n$. (1.13.A.1--A.5)
  \item Equal spacing is not a rule to memorize: a difference compares answers
        across \textbf{one step of the question}, so the steps must match.
  \item A degree-$n$ polynomial fits any $n+1$ points with distinct inputs, so
        ``it fits the data'' is not a rationale. (1.13.A.6)
\end{itemize} \\
\textbf{1.C} --- Construct a new function from a context, a transformation, or a
regression. \newline
\textbf{3.B} --- Apply numerical results in a given applied context. &
\begin{itemize}[leftmargin=1.2em,itemsep=2pt,topsep=0pt]
  \item Building from a restriction: name the question, write the constraint,
        solve it for the \textit{other} quantity, substitute into the one you
        want. (1.14.A.1--A.2)
  \item The clue in a context is how the answer \textbf{moves} as the question
        grows --- two lengths multiplied $\Rightarrow$ quadratic, three
        $\Rightarrow$ cubic, question underneath $\Rightarrow$ rational, rule
        changing at a threshold $\Rightarrow$ piecewise. (1.13.A.3, A.7,
        1.14.B.1)
  \item A model is used, not just written: evaluate it, and compute
        $\frac{f(b)-f(a)}{b-a}$ \textbf{with units} --- output units per input
        unit. (1.14.C.1)
\end{itemize} \\
\textbf{Synthesis} --- the choice of question is visible in every later answer. &
\begin{itemize}[leftmargin=1.2em,itemsep=2pt,topsep=0pt]
  \item Two students modeling one pen get different formulas, different
        domains, and different rate units --- and both are right.
  \item The domain is the set of questions the \textbf{context} allows, not
        what the algebra tolerates. Negative widths are not errors; they are
        questions the situation forbids. (1.13.B.1--B.4)
  \item Naming the assumption is part of the answer. A model that fits the
        numbers can still be wrong about the world.
\end{itemize} \\
\end{tabularx}
\end{skillbox}

\begin{skillbox}[{Vocabulary, Concepts, \& Theorems}]{greenbox}
\renewcommand{\arraystretch}{1.4}
\begin{tabularx}{\linewidth}{@{} >{\leavevmode\bfseries\color{plum}}p{0.28\linewidth} X @{}}
Function model        & A function built to stand in for a real situation, so the situation can be predicted and measured. \\
Question (input)      & The one quantity the model takes in. Choosing it is the first act of modeling; every other quantity must become an answer to it. \\
Constraint            & The equation the context forces on the quantities, such as ``the fencing totals 80 feet.'' It is what lets two variables collapse into one. \\
First differences     & The differences between consecutive outputs of \textbf{equally spaced} inputs. Constant first differences mean a linear model. \\
Second differences    & The differences of the first differences. Constant second differences mean a quadratic model --- the rate of change is itself changing at a constant rate. \\
Domain restriction    & The questions the \textit{context} allows, which is usually narrower than what the algebra allows. \\
Piecewise-defined function & A function given by different rules on different intervals; used when the constraint itself changes at a threshold. \\
Inverse proportion    & \TallMath{y = \frac{k}{x^n}} --- the question sits underneath, so as it grows the answer shrinks. Modeled by a rational function; Unit 3 takes this up in full. \\
Model assumption      & The simplifying claim about the world that has to be true for the model to be trusted. \\
\end{tabularx}
\end{skillbox}

\begin{skillbox}[Activate Prior Knowledge \& Spiral Review]{lilac}
    Please see attached warm-up (spiral review).
    Skills reviewed: computing differences down an output column and noticing
    when the \textit{input} row is not equally spaced (Lessons 1.2--1.3);
    solving one literal equation, $2w + \ell = 80$, for \textbf{each} of its two
    variables; and computing an average rate of change from a table with correct
    units (Lesson 1.3). Item 2 is the load-bearing one --- the two solves are
    Ana's model and Ben's model, built before either has a name. Item 1's second
    table returns in Part 2, where the class discovers it was linear all along.
\end{skillbox}

\begin{skillbox}[Hook]{lilac}
    \textbf{Two notebooks, one pen.} The garden club has exactly 80 feet of
    fencing for a rectangular pen; one whole side runs along the barn wall. Two
    members measured the \textit{same five trial pens}. Ana recorded
    width~$\to$~area ($5 \to 350$, $10 \to 600$, $15 \to 750$, $20 \to 800$,
    $25 \to 750$); Ben recorded length~$\to$~area ($30 \to 750$, $40 \to 800$,
    $50 \to 750$, $60 \to 600$, $70 \to 350$).
    \begin{enumerate}[itemsep=2pt]
        \item Put both tables on the board and ask, plainly: \textit{which
              member recorded the pen correctly?} Let the class argue. The
              answer is \textbf{both}, and nobody gets there in under a minute.
        \item Then the second vote, before any algebra: \textit{should the pen
              be a square?} Most classes say yes. Record it --- it will be
              wrong, and that is the point.
    \end{enumerate}
    \textit{Discussion prompt:} the two tables differ because the two members
    picked different \textbf{questions}. That single decision is the lesson.
\end{skillbox}

\begin{skillbox}[Lesson]{lilac}
\begin{multicols}{2}

\textbf{Part 1: Pick the Question} (6 min)

\begin{itemize}
  \item Name what changed: for six lessons students were handed a function.
        Today they are handed a \textbf{situation}, and a situation is not a
        function.
  \item Four numbers live in the pen: 80 (fixed), $w$, $\ell$, $A$. A function
        takes one question and returns one answer, so two decisions come before
        any algebra. Today the answer is settled --- the club wants area --- so
        only the question is open.
  \item Resolve the hook here: Ana's question was the width, Ben's the length,
        the pens were the same. Then name the 80 for what it is --- neither
        question nor answer, but the \textbf{constraint}. Write
        $2w + \ell = 80$ and leave it up all period.
\end{itemize}

\textbf{Part 2: When the Table Already Picked It} (10 min)

\begin{itemize}
  \item The top row of a table \textit{is} the question. Table A
        ($y = 5,8,11,14,17$): first differences all 3 $\Rightarrow$ linear
        (1.13.A.1--A.2).
  \item Derive the equal-spacing requirement instead of decreeing it: a
        difference compares two answers across \textbf{one step of the
        question}, so the steps have to match. Then run the warm-up's $g$ table
        ($x = 1, 2, 4$): differences 4 and 8, but \textit{per unit} both are 4
        --- $g$ was linear all along, and the raw differences lied.
  \item Ana's table: first differences $250, 150, 50, -50$ are not constant;
        the second differences are all $-100$ $\Rightarrow$ quadratic
        (1.13.A.4). Say the meaning: constant first differences mean the
        \textit{amount} of change repeats; constant second differences mean the
        \textit{rate} changes at a steady rate. Extend aloud to cubic
        (1.13.A.5).
  \item Close with the counting fact --- $n+1$ points fit a degree-$n$
        polynomial (1.13.A.6) --- and use it to motivate Part 3: a table alone
        cannot settle the model, so the context gets a vote.
\end{itemize}

\columnbreak

\textbf{Part 3: Reading the Context} (7 min)

\begin{itemize}
  \item Fill the clue table together. Every row is phrased as \textit{what
        happens to the answer as the question grows}, which is the only form of
        the clue worth memorizing.
  \item Then the sharpening, and it is the item most likely to move a test
        score: ``area'' does not mean quadratic. Length always twice the width
        gives $A(w) = 2w^2$; length \textbf{fixed} at 9 feet gives $A(w) = 9w$.
        Count the lengths that move with the question (1.13.A.3).
  \item Inverse proportion $\Rightarrow$ rational is \textbf{exposure only}:
        the question sits underneath, show $y = k/x^n$, name Unit 3
        (1.14.B.1). A rule that changes at a threshold $\Rightarrow$ piecewise
        (1.13.A.7).
\end{itemize}

\textbf{Part 4: Build It} (14 min)

\begin{itemize}
  \item The heart of the lesson. Four steps, said out loud every time: name the
        question, write the constraint, solve the constraint for the
        \textit{other} quantity, substitute into the one you want
        (1.14.A.1--A.2).
  \item Ana: $\ell = 80 - 2w \Rightarrow A(w) = w(80-2w) = 80w - 2w^2$. Check
        $A(20) = 800$ against her table. Then tie it back: Part 3 predicted
        quadratic because \textit{both} lengths move with $w$.
  \item Ben, same constraint solved for the other letter:
        $w = (80-\ell)/2 \Rightarrow A(\ell) = 40\ell - \ell^2/2$, and
        $A(40) = 800$. \textbf{Two formulas, one pen.} Say that sentence.
\end{itemize}

\textbf{Part 5: Use the Model} (9 min)

\begin{itemize}
  \item Evaluate $A(10) = A(30) = 600$. Average rate of change on $[10,20]$ is
        $+20$, on $[20,30]$ is $-20$ (1.14.C.1).
  \item \textbf{Units are not optional}, and here they do real work: square feet
        per foot of \textit{width} for Ana, per foot of \textit{length} for Ben
        (whose rates straddling the same pen, on $[30,40]$ and $[40,50]$, are
        $+5$ and $-5$). The units name the question.
  \item Sign change locates the maximum: $w = -b/(2a) = 20$, $\ell = 40$,
        $A = 800$. Return to the vote --- the best pen is $20 \times 40$,
        \textbf{not} a square. Ben lands on the same pen from the other end.
\end{itemize}

\textbf{Part 6: The Fine Print} (7 min)

\begin{itemize}
  \item Domain as a list of legal \textit{questions}: Ana $0 < w < 40$, Ben
        $0 < \ell < 80$ --- same pen, different domains (1.13.B.3). Range
        $0 < A \le 800$ (1.13.B.4).
  \item $A(-3) = -258$ is not bad arithmetic; it answers a question the
        situation forbids.
  \item Assumptions (1.13.B.1--B.2): all 80 ft used, barn at least 40 ft, pen
        exactly rectangular, no gate gap. If the barn is only 30 ft the
        constraint changes past a threshold --- \textbf{piecewise}, a callback
        to Part 3 and the homework extension.
\end{itemize}
\end{multicols}
\end{skillbox}

\begin{skillbox}[Active Monitoring]{redbox}
While students work on guided notes and practice, circulate and check:
\begin{itemize}
  \item \textbf{Common error:} taking differences when the inputs are not
        equally spaced. This is the lesson's signature error and it is silent
        --- the numbers still subtract. Ask ``what is the gap between your
        questions?'' before you ask what the differences are.
  \item \textbf{Common error:} stopping at the first differences, seeing that
        they are not constant, and concluding ``not a polynomial.'' Prompt:
        ``take the differences again.''
  \item \textbf{Common error:} in Part 4, writing $A = w\ell$ and stopping. A
        model must be a function of the \textbf{one} quantity you named as the
        question. Ask: ``what does the constraint let you replace?''
  \item \textbf{Common error:} $2w + \ell = 80$ becoming $\ell = 80 - w$ or
        $\ell = 40 - 2w$. Catch it early --- every later answer inherits it.
  \item \textbf{Common error:} treating ``area'' as automatically quadratic.
        The fixed-length rug is the counterexample; keep it on the board.
  \item \textbf{Common error:} the domain given as ``all real numbers'' or
        ``$w > 0$'' with no upper end. Push for the sentence: at $w = 40$ there
        is no fencing left for the length.
  \item \textbf{Common error:} an average rate of change reported as a bare
        number. No units, no credit --- and the units are where Ana's model and
        Ben's model become distinguishable.
\end{itemize}
\end{skillbox}

\begin{skillbox}[Group Work \& Differentiation]{redbox}
Students work on the \textbf{Group Activity: The Cardboard Bin} in leveled
tiers. Every tier uses the same sheet --- 24 in by 18 in, a square of side $x$
cut from each corner --- so groups can be re-sorted mid-period without new
setup.

\begin{itemize}
  \item \textbf{Tier R (Remediate):} Name the question in the club's table,
        check that the cuts are equally spaced, take differences down to the
        constant \textit{third} row, name cubic, and match six scenarios to
        their model types --- including the two rugs from Part 3.
  \item \textbf{Tier A (Apply):} Write each dimension of the bin in terms of
        the cut, build $V(x) = x(24-2x)(18-2x) = 4x^3 - 84x^2 + 432x$, check
        $V(3) = 648$ against the table, restrict to $0 < x < 9$, and find that
        the untested cut $x = 3.5$ beats every bin the club actually built.
  \item \textbf{Tier E (Extend):} Two average rates of change with units; then
        the same sheet with a \textit{different answer} --- the shipping
        footprint $B(x) = 4x^2 - 84x + 432$, quadratic where the volume was
        cubic; then a \textit{different question} --- fix the volume at
        648~in$^3$ and get $h(B) = 648/B$, an inverse proportion and a bridge
        to Unit 3.
\end{itemize}
\end{skillbox}

\begin{skillbox}[Individual Work \& Assessment]{redbox}
\textbf{Exit Ticket} (last 5--7 minutes, individual, collected):
\begin{enumerate}
  \item Name which row is the question, complete a difference table, and name
        the model type it points to.
  \item Build an area model from a 40-foot fencing constraint, state its
        domain, and write one assumption it makes.
  \item Two rugs --- one with length always three times the width, one with a
        fixed 8-foot length --- write each area model, name each type, and
        explain why only one is quadratic.
\end{enumerate}
\textbf{Assessment note:} Item 2 is the diagnostic for construction (1.14.A) ---
a student who writes $A(w) = w(40-w)$ has forgotten that two widths are fenced
and needs Part 4 retaught in tomorrow's warm-up window. \textbf{Item 3 is the
spine check and the fastest scan on the page:} a student who calls both rugs
quadratic is matching nouns, not tracking how the answer moves with the
question, and that is the habit Unit 2 will punish.
\end{skillbox}

\begin{skillbox}[Reinforcement \& Extension]{goldbox}
\textbf{Homework:} 6 problems covering difference tables (including a cubic),
matching contexts to model types, \textbf{building Ben's model} for the same pen
and comparing its domain to Ana's, using Ana's model (evaluate, average rate of
change with units, domain check), an inverse-proportion model, and two AP-style
multiple-choice items --- one on constant second differences, one on the
fixed-height flag whose area is linear.

\textbf{Extension (optional):} The barn wall turns out to be only 30 feet long,
so widths below 25 feet need fencing on part of the fourth side too. Students
find the threshold, write the two rules, check that they agree there, and
explain why the model is now piecewise.

\textbf{Preview:} Unit 1 ends here. Lesson 2.1 (Quadratic Functions Revisited)
opens Unit 2 with $A(w) = 80w - 2w^2$ still on the board --- vertex form, the
axis of symmetry, and why $-b/(2a)$ works instead of just that it does. Every
model chosen today was a polynomial; Unit 2 is where polynomials get studied on
their own terms.
\end{skillbox}

\begin{teachernote}[Warm-Up]
Five minutes, silent start, calculators fine. Item 1's second table is the only
place all period where students see inputs that are \textit{not} equally spaced
--- part (b) is the trap, and it is worth ten seconds of ``look at the $x$ row
first'' when you review. Do not resolve it here; Part 2 rehabilitates that table
and shows $g$ was linear after all, which is a much better moment than a
correction. Item 2 does the Part 4 algebra twice with no context attached, and
students will ask why they are solving the same equation for both letters ---
say ``you'll see'' and mean it. In Part 4, put the warm-up back on the board:
they built both of today's models before either had a name. Item 3 is Lesson 1.3
with units, and units are what Part 5 grades on.
\end{teachernote}

\begin{teachernote}[Guided Notes]
Part 1 is short but it is the lesson; do not let it become a vocabulary
recitation. The moment worth protecting is the class realizing that Ana and Ben
are \textit{both} right --- give it a full minute of argument before you
resolve it, because every later payoff (two formulas, two domains, two sets of
rate units) collects on that minute. Part 2's equal-spacing requirement must be
said before the first subtraction, and derive it rather than decree it: a
difference compares answers across one step of the question. Part 3's two rugs
are the highest-value thirty seconds on the page --- if a student can say why
the fixed-length rug is linear, they understand the whole selection half. Part 4
is where the period is won or lost: do all four steps out loud, and do not let
the class skip the substitution because ``it's obvious.'' If time runs short,
compress Part 6 to the two domains and one assumption, and let the homework
carry the range and the piecewise wrinkle --- but do not cut Ben's model out of
Part 4, because the domain comparison in Part 6 and homework problem 3 both
depend on it.
\end{teachernote}

\begin{teachernote}[Group Activity]
All three tiers run on the same 24-by-18 sheet, so a group that finishes Tier R
can move to Tier A without re-reading the scenario. Bring one real sheet of
cardboard and a pair of scissors if you can --- cutting the corners once, in
front of the class, makes ``the only thing you control is the cut'' land
physically, and it is the whole reason $x$ is the question. Tier R's third
difference row is the first time students carry differences past the second row
on real data; expect arithmetic slips and let groups check each other rather
than checking for them. Tier A's step 1 is the one to watch as you circulate: a
group that writes the base as $24 - x$ instead of $24 - 2x$ has forgotten that
both ends are cut, and everything downstream is wrong and quiet about it. Tier
A step 5 is the payoff --- the untested cut $x = 3.5$ beats all six bins the
club built, which is the cleanest argument in the lesson for why you want a
model and not a table. In Tier E, item 2 changes the \textit{answer} and item 3
changes the \textit{question}; that contrast is deliberate, and it is worth
naming aloud to any group that reaches it. If the class never reaches Tier E,
assign homework problem 3 as the substitute.
\end{teachernote}

\begin{teachernote}[Exit Ticket]
Collected, completion credit only. Sort by item 3 first --- it is four seconds
per paper and it separates ``matched the word \textit{area} to the word
\textit{quadratic}'' from ``counted the lengths that move.'' Then sort by item 2.
Accept $A(w) = 40w - 2w^2$ or $A(w) = w(40 - 2w)$; the factored form often
signals a student who understood the substitution rather than memorizing a
formula. A domain of $0 < w < 20$ is the target; $w > 0$ alone is half credit
and goes in the reteach pile. On item 2(d), any defensible assumption earns the
point --- the skill is that the student wrote one at all.
\end{teachernote}

\begin{teachernote}[Homework]
Problem 3 is the flagship and the one to scan first: it is Ben's model, and part
(d) --- why two correct models of one pen have different domains --- is the
single best predictor of who understood the day. Problem 4(b)'s units predict
who is ready for Unit 2's modeling work. Problem 5 is the only
inverse-proportion practice students get before Unit 3; expect the arithmetic on
$k = F d^2$ to be shakier than the concept, and grade the domain sentence, not
the decimal. On problem 6(a), distractor A is the student who never took the
second differences, C is the student who counted rows instead of differences,
and D confuses ``constant second differences'' with ``second differences of
zero.'' Problem 6(b) is item 3 of the exit ticket in multiple-choice clothing;
if a student misses both, reteach Part 3 rather than Part 2. The extension is
genuinely hard for this track --- assign it to students who finished Tier E, not
to everyone.
\end{teachernote}
\end{document}
